\begin{abstract}
Existing text-to-image metrics are largely calibrated for single-subject alignment or global aesthetic preference, and they degrade sharply when multiple personalized subjects must be generated in the same image. In this regime, the key failure modes are not merely poor global semantics, but missing identities, local appearance entanglement, and incorrect interaction assignment across subjects. We argue that current evaluation practice therefore suffers from a structural blind spot: a metric may return similar scores for obviously different generations once the number of requested subjects becomes large. To address this gap, we propose PrismBench, a benchmark for multi-subject personalized generation built from paired generations under shared prompts and shared reference identities, and LENS, a reference-conditioned diagnostic evaluator trained to score pairwise preference while predicting three image-level error dimensions: Existence, Appearance, and Interaction. PrismBench combines a large silver training set labeled by dual multimodal teachers with a high-quality golden set labeled by double-blind human annotators. LENS is optimized with a dual-head objective that unifies pairwise ranking and diagnostic supervision. Although the experimental section is left for completion after large-scale data collection, the benchmark design, labeling protocol, and training formulation together establish a practical and testable route toward restoring score separability in high-subject-count generation.
\end{abstract}
